\clearpage
\section{多标签：逐个标签正确，还是整组正确？}
\label{sec:eval-multilabel}
一条记录可以同时有多个标签，例如一张图片既有“猫”也有“室内”。
用$y_{ik},\hat y_{ik}\in\{0,1\}$表示第$i$条记录、第$k$个标签的真实与预测状态，
共$n$条记录、$K$个标签。这与每条记录只选一个类别的多分类不同。

\begin{figure}[H]
\centering\begin{tikzpicture}[x=12mm,y=8mm]
\node[lab] at (1.5,4) {真实标签};\node[lab] at (7.5,4) {预测标签};
\foreach \x/\name in {0/A,1/B,2/C,3/D}{
 \node[lab] at (\x,3) {\name};\node[lab] at (\x+6,3) {\name};}
\foreach \y in {0,1,2}{\node[lab,anchor=east] at (-.6,\y) {样本\pgfmathparse{int(3-\y)}\pgfmathresult};}
\foreach \x/\y/\v in {0/2/1,1/2/0,2/2/1,3/2/0,0/1/0,1/1/1,2/1/0,3/1/0,0/0/1,1/0/1,2/0/0,3/0/0}{
 \node[draw=Teal,fill=Pale,minimum width=11mm,minimum height=7mm,font=\small] at (\x,\y) {\v};}
\foreach \x/\y/\v in {6/2/1,7/2/0,9/2/0,6/1/0,7/1/1,8/1/0,6/0/1,7/0/1,8/0/0,9/0/0}{
 \node[draw=Teal,fill=Pale,minimum width=11mm,minimum height=7mm,font=\small] at (\x,\y) {\v};}
\foreach \x/\y/\v in {8/2/0,9/1/1}{
 \node[draw=Amber,fill=Amber!15,minimum width=11mm,minimum height=7mm,font=\small] at (\x,\y) {\v};}
\draw[flow] (3.8,1)--(5.1,1);
\end{tikzpicture}

\caption{三个样本、四个标签；橙色格为两处错误。每一行是一个标签集合。}
\end{figure}
\readfigure{第1条漏掉标签C，第2条多报标签D，第3条完全正确。
12个二元位置错2个，但只有1条记录的整组标签完全匹配。}

\begin{align}
 L_{\mathrm{Hamming}}&=\frac1{nK}\sum_{i=1}^n\sum_{k=1}^K
      \mathbf1[y_{ik}\ne\hat y_{ik}],\\
 \mathrm{SubsetAccuracy}&=\frac1n\sum_{i=1}^n
      \mathbf1[\vct y_i=\hat{\vct y}_i].
\end{align}
本例Hamming损失为$2/12=1/6$，越小越好；整组匹配准确率为$1/3$，越大越好。
整组0--1损失为$1-\mathrm{SubsetAccuracy}=2/3$，不是Hamming损失。
这些定义不要求不同标签在统计上独立。

\subsection{F1也要说明在哪个方向汇总}
各标签先用各自计数求F1，再等权平均，是Macro-F1；先将所有标签的计数相加再求F1，是Micro-F1。
也可先对每条记录的标签集合求F1，再对记录平均，称Samples-F1。
本例总TP=4、FP=1、FN=1，Micro-F1为$8/10=0.8$；
各标签F1为$1,1,0,0$，Macro-F1为0.5；Samples-F1为$(2/3+2/3+1)/3=7/9$。

\begin{pitfall}[很多负标签会让Hamming看起来很小]
若100个标签中每条记录恰有1个正标签，全部预测0的Hamming损失只有0.01，
整组匹配准确率和召回率却为0。应结合正标签表现、各标签支持数与实际错误代价。
空集合导致某些F1分母为0时，另行说明软件约定。
\end{pitfall}

输出通常使用每标签Sigmoid和二元损失；各标签阈值在验证流程中选择。
损失按标签求和还是平均也需写明，不能直接比较不同缩放下的数值。
